\documentclass[a4paper,10pt]{article}
\usepackage[margin=2cm]{geometry}
\usepackage{listings}
\usepackage{xcolor}
\usepackage{fancyhdr}
\usepackage{ctex}
\usepackage{DejaVuSansMono}

\pagestyle{fancy}
\fancyhf{}
\fancyhead[C]{\small 审计智鉴流式数据处理系统 V1.0}
\fancyfoot[C]{\thepage/60}
\renewcommand{\headrulewidth}{0.4pt}

% 代码样式设置：纯黑白，无任何彩色
\lstset{
    language=Python,
    basicstyle=\ttfamily\footnotesize,
    numbers=left, numberstyle=\tiny,
    numbersep=8pt, frame=none, breaklines=true,
    breakatwhitespace=true, showstringspaces=false,
    tabsize=4, aboveskip=0pt, belowskip=0pt,
    lineskip=1pt,
    morekeywords={async,await,True,False,None},
    literate={_}{\_}1
}

\begin{document}
\setcounter{page}{31}

% ==================== frontend/app.py: render_header() ====================
\begin{lstlisting}
def render_header():
    """渲染顶部水平导航栏 - 替代侧边栏"""
    if "current_page" not in st.session_state:
        st.session_state.current_page = "home"
    st.markdown("""
    <div style="position:fixed;top:0;left:0;right:0;
        z-index:10000;height:3px;background:linear-gradient(
        90deg,#2563EB,#3B82F6,#60A5FA,#3B82F6,#2563EB);">
    </div>
    """, unsafe_allow_html=True)
    col_logo, col_nav, col_user = st.columns([1.2, 5, 1.3])
    with col_logo:
        st.markdown("""
        <div style="display:flex;align-items:center;
            gap:10px;cursor:pointer;">
            <div style="width:36px;height:36px;
                border-radius:10px;background:linear-gradient(
                135deg,#2563EB,#1D4ED8);display:flex;
                align-items:center;justify-content:center;
                color:#FFFFFF;font-weight:800;font-size:1rem;">
                A</div>
            <span style="font-size:1.25rem;font-weight:800;
                color:#0F172A;letter-spacing:-0.5px;">
                Audit Mind</span>
        </div>
        """, unsafe_allow_html=True)
    with col_nav:
        if st.session_state.logged_in:
            pages = [
                ("首页", "home"), ("财务助手", "fs"),
                ("舞弊检测", "detect"), ("AI 问答", "qa"),
                ("我的检测", "history"), ("报告", "reports"),
                ("会员", "membership"),
            ]
        else:
            pages = [
                ("首页", "home"), ("AI 问答", "qa"),
                ("价格", "pricing"), ("案例", "cases"),
            ]
        current = st.session_state.current_page
        nav_cols = st.columns(len(pages))
        for i, (label, key) in enumerate(pages):
            with nav_cols[i]:
                active = current == key
                if active:
                    st.markdown(f"""
                    <div style="text-align:center;padding:6px 2px;">
                        <span style="font-size:0.95rem;
                            font-weight:700;color:#2563EB;">
                            {label}</span>
                        <div style="width:20px;height:3px;
                            background:linear-gradient(90deg,
                            #2563EB,#3B82F6);border-radius:2px;
                            margin:4px auto 0;"></div>
                    </div>
                    """, unsafe_allow_html=True)
                else:
                    if st.button(label, key=f"nav_{key}",
                            use_container_width=True):
                        st.session_state.current_page = key
                        st.rerun()
    with col_user:
        if st.session_state.logged_in:
            user_cols = st.columns([2, 1])
            with user_cols[0]:
                username = st.session_state.user_info.get(
                    'username', '用户')
                st.markdown(f"""
                <div style="display:flex;align-items:center;
                    justify-content:flex-end;gap:8px;
                    padding-top:4px;">
                    <div style="width:32px;height:32px;
                        border-radius:50%;background:linear-gradient(
                        135deg,#2563EB,#3B82F6);display:flex;
                        align-items:center;justify-content:center;
                        color:#FFFFFF;font-weight:700;
                        font-size:0.85rem;">
                        {username[0].upper()}</div>
                    <span style="font-size:0.9rem;font-weight:600;
                        color:#0F172A;">{username}</span>
                </div>
                """, unsafe_allow_html=True)
            with user_cols[1]:
                if st.button("退出", key="header_logout",
                        use_container_width=True):
                    AuthManager.clear_auth()
                    st.session_state.logged_in = False
                    st.session_state.token = None
                    st.session_state.user_info = None
                    st.rerun()
        else:
            if st.button("登录 / 注册", key="header_login",
                    use_container_width=True, type="primary"):
                st.session_state.show_login_modal = True
                st.rerun()
\end{lstlisting}

% ==================== frontend/app.py: render_home() ====================
\begin{lstlisting}
def render_home():
    """渲染首页 - Cinematic Full-Width Design"""
    hero_html = '''
    <div style="width:100%;min-height:580px;position:relative;
        overflow:hidden;background:linear-gradient(135deg,
        #0F172A 0%,#1E293B 40%,#0F172A 100%);
        display:flex;align-items:center;justify-content:center;">
        <div style="position:absolute;top:0;left:0;right:0;
            bottom:0;opacity:0.15;background-image:linear-gradient(
            rgba(59,130,246,0.3) 1px,transparent 1px),
            linear-gradient(90deg,rgba(59,130,246,0.3) 1px,
            transparent 1px);background-size:60px 60px;"></div>
        <div style="position:relative;z-index:2;text-align:center;
            max-width:900px;padding:0 4rem;">
            <div style="display:inline-flex;align-items:center;
                gap:10px;padding:10px 24px;border-radius:100px;
                font-size:0.9rem;font-weight:600;background:
                rgba(37,99,235,0.15);border:1px solid
                rgba(96,165,250,0.3);color:#60A5FA;
                margin-bottom:2rem;">
                <span style="width:8px;height:8px;background:#3B82F6;
                    border-radius:50%;display:inline-block;
                    animation:pulse 2s infinite;"></span>
                AI-Powered Financial Audit Intelligence
            </div>
            <h1 style="font-size:4.5rem;font-weight:800;
                letter-spacing:-3px;line-height:1.05;
                margin-bottom:1.25rem;color:#FFFFFF;">
                智能识别财务舞弊
            </h1>
            <h2 style="font-size:2.4rem;font-weight:500;
                color:#94A3B8;letter-spacing:-1px;
                margin-bottom:1.5rem;">
                守护每一笔资金的安全
            </h2>
            <p style="font-size:1.25rem;color:#94A3B8;
                line-height:1.8;max-width:620px;
                margin:0 auto 2.5rem;">
                基于生成式 AI 的上市公司财务舞弊识别系统
                <br>双模分析 · SHAP 可解释 · 风险标签可视化
            </p>
        </div>
    </div>
    '''
    st.components.v1.html(hero_html, height=580, scrolling=False)
    stats_html = '''
    <div style="width:100%;padding:64px 0;background:
        linear-gradient(90deg,#2563EB,#3B82F6,#1D4ED8);">
        <div style="max-width:1200px;margin:0 auto;
            padding:0 4rem;display:grid;
            grid-template-columns:repeat(4,1fr);gap:2rem;">
            <div style="text-align:center;">
                <div style="font-size:3.5rem;font-weight:800;
                    color:#FFFFFF;">92%</div>
                <div style="font-size:1.05rem;
                    color:rgba(255,255,255,0.85);">
                    AI 识别准确率</div>
            </div>
            <div style="text-align:center;">
                <div style="font-size:3.5rem;font-weight:800;
                    color:#FFFFFF;">12,847</div>
                <div style="font-size:1.05rem;
                    color:rgba(255,255,255,0.85);">
                    已分析财报</div>
            </div>
            <div style="text-align:center;">
                <div style="font-size:3.5rem;font-weight:800;
                    color:#FFFFFF;">4,562</div>
                <div style="font-size:1.05rem;
                    color:rgba(255,255,255,0.85);">
                    覆盖 A 股企业</div>
            </div>
            <div style="text-align:center;">
                <div style="font-size:3.5rem;font-weight:800;
                    color:#FFFFFF;">7</div>
                <div style="font-size:1.05rem;
                    color:rgba(255,255,255,0.85);">
                    风险分析维度</div>
            </div>
        </div>
    </div>
    '''
    st.components.v1.html(stats_html, height=200, scrolling=False)
    c1, c2, c3 = st.columns([1, 1, 4])
    with c1:
        if st.button("开始检测", type="primary",
                use_container_width=True, key="hero_cta_detect"):
            if st.session_state.logged_in:
                st.session_state.current_page = "detect"
                st.rerun()
            else:
                st.session_state.show_login_modal = True
                st.rerun()
    with c2:
        if st.button("查看价格", use_container_width=True,
                key="hero_cta_more"):
            st.session_state.current_page = "pricing"
            st.rerun()
\end{lstlisting}

% ==================== frontend/app.py: render_detection() ====================
\begin{lstlisting}
def render_detection():
    """渲染舞弊检测页面 - Premium Design"""
    st.markdown("""
    <div style="margin:-1rem -1rem 2rem -1rem;
        padding:2rem 1.5rem;background:#F8FAFC;
        border-bottom:1px solid #E2E8F0;
        border-radius:20px;position:relative;
        overflow:hidden;">
        <div style="position:absolute;top:-50px;right:-50px;
            width:200px;height:200px;background:
            rgba(37,99,235,0.08);border-radius:50%;"></div>
        <div style="position:relative;z-index:1;">
            <h2 style="color:#0F172A;font-size:1.8rem;
                font-weight:700;margin-bottom:0.5rem;">
                舞弊检测</h2>
            <p style="color:#475569;font-size:1rem;margin:0;">
                上传财务数据，AI 智能识别潜在舞弊风险</p>
        </div>
    </div>
    """, unsafe_allow_html=True)
    if not st.session_state.logged_in:
        st.warning("请先登录以使用检测功能")
        return
    if 'upload_years' not in st.session_state:
        st.session_state.upload_years = [2023]
    if 'uploaded_files_data' not in st.session_state:
        st.session_state.uploaded_files_data = {}
    if 'parsed_results' not in st.session_state:
        st.session_state.parsed_results = None
    tab1, tab2, tab3 = st.tabs(
        [" 文件上传", " 内置案例库", "手动录入"])
    with tab1:
        st.subheader("上传财务文件进行检测")
        st.info("**上传说明**：您可以将多年度财务数据整理在
            一个Excel/CSV文件中上传")
        col1, col2 = st.columns(2)
        with col1:
            company_name = st.text_input("企业名称*",
                key="upload_company_name",
                placeholder="例如：贵州茅台")
        with col2:
            stock_code = st.text_input("证券代码",
                key="upload_stock_code",
                placeholder="例如：600519")
        st.subheader("设置数据年份区间")
        col_start, col_end = st.columns(2)
        with col_start:
            start_year = st.number_input("起始年份",
                min_value=2000, max_value=2025,
                value=2020, key="start_year")
        with col_end:
            end_year = st.number_input("结束年份",
                min_value=2000, max_value=2025,
                value=2023, key="end_year")
        if start_year > end_year:
            st.error("起始年份不能大于结束年份")
        uploaded_file = st.file_uploader(
            "上传财务数据文件",
            type=['xlsx', 'xls', 'csv', 'txt'],
            key="financial_file_uploader")
        mdna_files = st.file_uploader(
            "上传MD&A文本文件（可选，可多选）",
            type=['txt', 'docx', 'doc', 'pdf'],
            key="mdna_file_uploader",
            accept_multiple_files=True)
        if uploaded_file:
            if st.button("解析文件内容", type="secondary",
                    use_container_width=True):
                with st.spinner("正在解析文件，请稍候..."):
                    try:
                        file_content = uploaded_file.getvalue()
                        if uploaded_file.name.endswith(
                                ('.xlsx', '.xls')):
                            import pandas as pd
                            df = pd.read_excel(
                                io.BytesIO(file_content))
                        elif uploaded_file.name.endswith('.csv'):
                            import pandas as pd
                            df = pd.read_csv(
                                io.BytesIO(file_content))
                        year_col = None
                        for col in df.columns:
                            if str(col).lower() in \
                                    ['year', '年份', '年度']:
                                year_col = col
                                break
                        results = []
                        if year_col:
                            for year_val, year_df in \
                                    df.groupby(year_col):
                                if start_year <= int(year_val) \
                                        <= end_year:
                                    financial_data = {}
                                    for col in year_df.columns:
                                        if col != year_col:
                                            val = year_df[col].iloc[0]
                                            try:
                                                financial_data[
                                                    str(col)] = float(val)
                                            except:
                                                pass
                                    results.append({
                                        "year": int(year_val),
                                        "financial_data": financial_data,
                                        "parsed_success": True
                                    })
                        else:
                            financial_data = {}
                            for col in df.columns:
                                try:
                                    val = df[col].iloc[0]
                                    financial_data[str(col)] = float(val)
                                except:
                                    pass
                            results.append({
                                "year": start_year,
                                "financial_data": financial_data,
                                "parsed_success": True
                            })
                        mdna_text = ""
                        if mdna_files:
                            mdna_parts = []
                            for i, mdna_f in enumerate(mdna_files):
                                content = mdna_f.getvalue().decode(
                                    'utf-8', errors='ignore')
                                if content.strip():
                                    mdna_parts.append(
                                        f"【MD&A文件{i+1}: {mdna_f.name}】\n{content}")
                            mdna_text = "\n\n---\n\n".join(mdna_parts)
                        st.session_state.parsed_results = {
                            "type": "multi_year",
                            "results": results,
                            "year_range": f"{start_year}-{end_year}",
                            "mdna_text": mdna_text
                        }
                        st.success(f"成功解析 {len(results)} 个年份的数据")
                        st.rerun()
                    except Exception as e:
                        st.error(f"解析失败：{str(e)}")
\end{lstlisting}

% ==================== frontend/app.py: render_detection_result() ====================
\begin{lstlisting}
def render_detection_result(result, show_divider=True):
    """渲染检测结果 - Premium Design"""
    year = result.get('year', '')
    fraud_prob = result.get("fraud_probability", 0)
    risk_level = result.get("risk_level", "low")
    risk_score = result.get("risk_score", 0)
    risk_colors = {"high": "#EF4444", "medium": "#F59E0B",
        "low": "#10B981"}
    risk_bg = {"high": "rgba(239,68,68,0.1)",
        "medium": "rgba(245,158,11,0.1)",
        "low": "rgba(16,185,129,0.1)"}
    rc = risk_colors.get(risk_level, "#6B8294")
    rbg = risk_bg.get(risk_level, "rgba(107,130,148,0.1)")
    st.markdown(f"""
    <div style="margin:0 0 2rem 0;padding:2rem;
        border-radius:24px;background:linear-gradient(
        135deg,#F8FAFC,#EFF6FF);border:1px solid {rc}33;">
        <div style="display:flex;align-items:center;
            gap:12px;margin-bottom:1rem;">
            <span style="display:inline-flex;padding:6px 16px;
                border-radius:100px;font-size:0.85rem;
                font-weight:700;background:{rbg};
                color:{rc};border:1px solid {rc}44;">
                {show_risk_level_badge(risk_level)}</span>
            <span style="font-size:1.5rem;font-weight:700;">
                {year}年度 智能检测报告</span>
        </div>
        <div style="display:grid;
            grid-template-columns:repeat(4,1fr);gap:1rem;">
            <div style="text-align:center;padding:1rem;
                border-radius:16px;">
                <div style="font-size:0.8rem;opacity:0.6;
                    margin-bottom:4px;">舞弊概率</div>
                <div style="font-size:1.8rem;font-weight:800;
                    color:{rc};">{fraud_prob:.1%}</div>
            </div>
            <div style="text-align:center;padding:1rem;
                border-radius:16px;">
                <div style="font-size:0.8rem;opacity:0.6;
                    margin-bottom:4px;">风险评分</div>
                <div style="font-size:1.8rem;font-weight:800;">
                    {risk_score:.1f}</div>
            </div>
            <div style="text-align:center;padding:1rem;
                border-radius:16px;">
                <div style="font-size:0.8rem;opacity:0.6;
                    margin-bottom:4px;">风险标签</div>
                <div style="font-size:1.8rem;font-weight:800;">
                    {len(result.get('risk_labels', []))}</div>
            </div>
            <div style="text-align:center;padding:1rem;
                border-radius:16px;">
                <div style="font-size:0.8rem;opacity:0.6;
                    margin-bottom:4px;">置信度</div>
                <div style="font-size:1.8rem;font-weight:800;">
                    {(1-fraud_prob)*100:.0f}%</div>
            </div>
        </div>
    </div>
    """, unsafe_allow_html=True)
    col1, col2 = st.columns([1, 1])
    with col1:
        fig = create_fraud_probability_gauge(fraud_prob)
        st.plotly_chart(fig, use_container_width=True)
    with col2:
        st.metric("风险标签数",
            f"{len(result.get('risk_labels', []))}个")
    st.divider()
    col_chart1, col_chart2 = st.columns([2, 1])
    with col_chart1:
        st.subheader("AI风险特征雷达图")
        ai_scores = result.get("ai_feature_scores", {})
        if ai_scores:
            fig = create_ai_radar_chart(ai_scores)
            if fig:
                st.plotly_chart(fig, use_container_width=True)
    with col_chart2:
        st.subheader("SHAP特征重要性")
        shap_features = result.get("shap_features", {})
        if shap_features:
            sorted_features = sorted(shap_features.items(),
                key=lambda x: abs(x[1]), reverse=True)[:5]
            for feature, importance in sorted_features:
                abs_importance = abs(importance)
                display_value = min(abs_importance * 2, 1.0)
                feature_names = {
                    "CON_SEM_AI": "语义矛盾度",
                    "COV_RISK_AI": "风险披露完整性",
                    "TONE_ABN_AI": "异常乐观语调",
                    "FIT_TD_AI": "文本-数据一致性",
                    "HIDE_REL_AI": "关联隐藏指数",
                    "DEN_ABN_AI": "信息密度异常",
                    "STR_EVA_AI": "回避表述强度"
                }
                feature_name = feature_names.get(feature, feature)
                st.progress(display_value,
                    text=f"{feature_name}: {importance:+.4f}")
\end{lstlisting}

% ==================== frontend/app.py: render_qa() ====================
\begin{lstlisting}
def render_qa():
    """渲染 AI 问答页面 - 支持流式输出"""
    _render_page_header("AI 智能问答",
        "财务舞弊领域专业问答助手，7x24 随时解答")
    if not st.session_state.logged_in:
        st.info("AI 问答功能仅对登录用户开放")
        st.divider()
        render_login_register()
        return
    if not st.session_state.chat_history:
        st.markdown('''
        <div class="glass-card" style="text-align:center;
            padding:3rem 2rem;margin-bottom:1.5rem;">
            <h3 style="font-size:1.25rem;font-weight:700;
                margin-bottom:0.5rem;">
                财务舞弊智能问答助手</h3>
            <p style="font-size:0.95rem;opacity:0.7;
                line-height:1.6;max-width:500px;margin:0 auto;">
                我可以帮您解答财务舞弊识别、审计方法、
                案例分析等专业问题。
                <br>支持流式输出，即问即答。
            </p>
        </div>
        ''', unsafe_allow_html=True)
    if 'streaming_answer' not in st.session_state:
        st.session_state.streaming_answer = ""
    if 'is_streaming' not in st.session_state:
        st.session_state.is_streaming = False
    suggestions = make_api_request("/qa/suggestions")
    with st.sidebar:
        st.subheader("推荐问题")
        if suggestions:
            for cat in suggestions:
                with st.expander(cat["category"]):
                    for q in cat["questions"]:
                        if st.button(q, key=f"sug_{q[:20]}"):
                            st.session_state.selected_question = q
    for msg in st.session_state.chat_history:
        with st.chat_message(msg["role"]):
            st.markdown(msg["content"])
    if prompt := st.chat_input("请输入您的问题..."):
        st.session_state.chat_history.append(
            {"role": "user", "content": prompt})
        with st.chat_message("user"):
            st.markdown(prompt)
        with st.chat_message("assistant"):
            thinking_placeholder = st.empty()
            thinking_placeholder.markdown("""
            <div style="display:flex;align-items:center;
                gap:10px;padding:4px 0;">
                <div style="display:flex;gap:4px;
                    align-items:center;">
                    <span style="display:inline-block;
                        width:7px;height:7px;border-radius:50%;
                        background:#3B82F6;
                        animation:thinkingBounce 1.4s infinite
                        ease-in-out both;"></span>
                    <span style="display:inline-block;
                        width:7px;height:7px;border-radius:50%;
                        background:#3B82F6;
                        animation:thinkingBounce 1.4s infinite
                        ease-in-out both;
                        animation-delay:0.16s;"></span>
                    <span style="display:inline-block;
                        width:7px;height:7px;border-radius:50%;
                        background:#3B82F6;
                        animation:thinkingBounce 1.4s infinite
                        ease-in-out both;
                        animation-delay:0.32s;"></span>
                </div>
                <span style="font-size:0.85rem;color:#64748B;">
                    AI 正在思考中...</span>
            </div>
            """, unsafe_allow_html=True)
            try:
                import requests
                url = f"{API_BASE_URL}/qa/ask-stream"
                req_headers = {
                    "Authorization": f"Bearer {st.session_state.token}",
                    "Content-Type": "application/json"
                }
                req_data = {"question": prompt}
                response = requests.post(url, json=req_data,
                    headers=req_headers, stream=True, timeout=120)
                if response.status_code == 200:
                    def _stream_generator():
                        for line in response.iter_lines(
                                decode_unicode=True):
                            if not line:
                                continue
                            if line.startswith("data: "):
                                data_str = line[6:]
                                if data_str == "[DONE]":
                                    break
                                try:
                                    event_data = json.loads(data_str)
                                    content = event_data.get("content")
                                    if content and \
                                            isinstance(content, str):
                                        yield content
                                    if event_data.get("error"):
                                        break
                                except json.JSONDecodeError:
                                    continue
                    full_answer = st.write_stream(_stream_generator)
                    thinking_placeholder.empty()
                    st.session_state.chat_history.append(
                        {"role": "assistant", "content": full_answer})
                else:
                    thinking_placeholder.empty()
                    result = make_api_request("/qa/ask",
                        method="POST", data={"question": prompt})
                    if result and "answer" in result:
                        answer = result["answer"]
                        st.markdown(answer)
                        st.session_state.chat_history.append(
                            {"role": "assistant", "content": answer})
                    else:
                        st.error("回答失败，请稍后重试")
            except Exception as e:
                thinking_placeholder.empty()
                result = make_api_request("/qa/ask",
                    method="POST", data={"question": prompt})
                if result and "answer" in result:
                    answer = result["answer"]
                    st.markdown(answer)
                    st.session_state.chat_history.append(
                        {"role": "assistant", "content": answer})
                else:
                    st.error(f"回答失败: {str(e)}")
\end{lstlisting}

% ==================== frontend/app.py: render_tax_module() ====================
\begin{lstlisting}
def render_tax_module():
    """渲染税务中心模块 - 包含税务测算、简易报税、财税咨询"""
    sub_tab1, sub_tab2, sub_tab3 = st.tabs(
        ["税务测算", "简易报税", "财税咨询"])
    with sub_tab1:
        st.markdown("""
        <div style="text-align:center;padding:12px 0 16px 0;">
            <h3 style="font-size:1.2rem;font-weight:700;
                margin-bottom:4px;color:#0F172A;">
                智能税务测算</h3>
            <p style="color:#64748B;font-size:0.85rem;margin:0;">
                按最新税法计算个税、企业税，对比最优纳税方案</p>
        </div>
        """, unsafe_allow_html=True)
        calc_tab1, calc_tab2 = st.tabs(["个税测算", "企业税测算"])
        with calc_tab1:
            st.markdown("**基本信息**")
            col1, col2 = st.columns(2)
            with col1:
                monthly_salary = st.number_input("税前月薪（元）",
                    min_value=0, value=15000, step=500,
                    key="tax_salary")
                annual_bonus = st.number_input("年终奖（元）",
                    min_value=0, value=30000, step=1000,
                    key="tax_bonus")
            with col2:
                social_insurance = st.number_input(
                    "五险一金/月（元）", min_value=0,
                    value=3000, step=100, key="tax_si")
                special_deduction = st.number_input(
                    "专项附加扣除/月（元）", min_value=0,
                    value=2000, step=100, key="tax_sd")
            if st.button("计算个税", type="primary",
                    use_container_width=True, key="calc_pit"):
                annual_income = monthly_salary * 12
                annual_deduction = (
                    60000 + social_insurance * 12
                    + special_deduction * 12
                )
                taxable_income = max(0,
                    annual_income - annual_deduction)
                brackets = [
                    (0, 36000, 0.03, 0),
                    (36000, 144000, 0.10, 2520),
                    (144000, 300000, 0.20, 16920),
                    (300000, 420000, 0.25, 31920),
                    (420000, 660000, 0.30, 52920),
                    (660000, 960000, 0.35, 85920),
                    (960000, float('inf'), 0.45, 181920),
                ]
                tax = 0
                rate = 0
                quick_deduction = 0
                for low, high, r, qd in brackets:
                    if taxable_income > low:
                        rate = r
                        quick_deduction = qd
                if taxable_income > 0:
                    tax = taxable_income * rate - quick_deduction
                after_tax = annual_income - tax
                st.metric("年收入", f"¥{annual_income:,.0f}")
                st.metric("年扣除额", f"¥{annual_deduction:,.0f}")
                st.metric("适用税率", f"{rate*100:.0f}%")
                st.metric("年个税", f"¥{tax:,.0f}")
                st.metric("税后年收入", f"¥{after_tax:,.0f}")
        with calc_tab2:
            col1, col2 = st.columns(2)
            with col1:
                revenue = st.number_input("营业收入（元/年）",
                    min_value=0, value=500000,
                    step=10000, key="biz_revenue")
                cost = st.number_input("营业成本（元/年）",
                    min_value=0, value=300000,
                    step=10000, key="biz_cost")
            with col2:
                taxpayer_type = st.selectbox("纳税人类型",
                    ["小规模纳税人", "一般纳税人"],
                    key="biz_type")
                industry = st.selectbox("行业类型",
                    ["服务业（6%）", "交通运输/建筑（9%）",
                    "货物销售（13%）"], key="biz_industry")
            if st.button("计算企业税", type="primary",
                    use_container_width=True, key="calc_biz"):
                profit = revenue - cost
                if taxpayer_type == "小规模纳税人":
                    vat_rate = 0.03
                    vat = revenue * vat_rate
                else:
                    vat_rates = {
                        "服务业（6%）": 0.06,
                        "交通运输/建筑（9%）": 0.09,
                        "货物销售（13%）": 0.13
                    }
                    vat_rate = vat_rates.get(industry, 0.06)
                    vat = revenue * vat_rate
                if profit <= 0:
                    income_tax = 0
                elif profit <= 3000000:
                    income_tax = profit * 0.05
                else:
                    income_tax = profit * 0.25
                surcharge = vat * 0.12
                if taxpayer_type == "小规模纳税人":
                    surcharge = vat * 0.06
                total_tax = vat + income_tax + surcharge
                net_profit = profit - total_tax
                st.metric("增值税", f"¥{vat:,.0f}")
                st.metric("企业所得税", f"¥{income_tax:,.0f}")
                st.metric("附加税", f"¥{surcharge:,.0f}")
                st.metric("合计应缴税额", f"¥{total_tax:,.0f}")
                st.metric("税后净利润", f"¥{net_profit:,.0f}")
\end{lstlisting}

% ==================== frontend/app.py: main() ====================
\begin{lstlisting}
def main():
    """主程序"""
    cookie_manager = AuthManager.get_cookie_manager_instance()
    if not st.session_state.auth_initialized:
        AuthManager.try_auto_login()
        st.session_state.auth_initialized = True
    render_header()
    if st.session_state.get('show_login_modal', False) \
            and not st.session_state.logged_in:
        render_login_modal()
    page = st.session_state.current_page
    non_fullwidth_pages = ["fs", "detect", "qa", "history",
        "reports", "membership", "settings", "pricing", "cases"]
    needs_wrap = page in non_fullwidth_pages
    if needs_wrap:
        st.markdown('<div class="content-wrap" style="
            padding-top:2rem;padding-bottom:3rem;">',
            unsafe_allow_html=True)
    if st.session_state.logged_in:
        if page == "home":
            render_home()
        elif page == "fs":
            render_financial_assistant()
        elif page == "detect":
            render_detection()
        elif page == "qa":
            render_qa()
        elif page == "history":
            render_my_detections()
        elif page == "reports":
            render_report_management()
        elif page == "membership":
            render_membership()
        elif page == "settings":
            render_account_settings()
    else:
        if page == "home":
            render_home()
        elif page == "qa":
            render_qa()
        elif page == "pricing":
            st.markdown("""
            <div style="margin:-2rem -2rem 2rem -2rem;
                padding:3rem 2rem;background:linear-gradient(
                135deg,#0F172A,#1E293B);">
                <div style="position:relative;z-index:1;
                    max-width:1280px;margin:0 auto;">
                    <h1 style="color:#FFFFFF;font-size:2.5rem;
                        font-weight:800;margin-bottom:0.5rem;">
                        价格中心</h1>
                    <p style="color:#94A3B8;font-size:1.05rem;">
                        选择适合您的方案，解锁全部智能审计能力</p>
                </div>
            </div>
            """, unsafe_allow_html=True)
            render_membership()
        elif page == "cases":
            render_case_center()
    if needs_wrap:
        st.markdown('</div>', unsafe_allow_html=True)

if __name__ == "__main__":
    main()
\end{lstlisting}

\end{document}
